\section{Абсолютная и условная сходимость}

Пусть теперь члены ряда $\displaystyle\sum_{k=1}^\infty a_k$ могут иметь любой знак.

\begin{definition}
Ряд $\sum a_k$ \emph{сходится абсолютно}, если сходится ряд из модулей $\sum\abs{a_k}$.
\end{definition}

\begin{theorem}[о связи сходимостей]
Абсолютно сходящийся ряд сходится.
\end{theorem}

\begin{proof}
Пусть $\sum\abs{a_k}$ сходится. По критерию Коши $\forall\eps>0\ \exists N\ \forall n>N\ \forall
p$: $\abs{a_{n+1}}+\dots+\abs{a_{n+p}}<\eps$. Тогда по неравенству треугольника
$\abs{a_{n+1}+\dots+a_{n+p}}\le\abs{a_{n+1}}+\dots+\abs{a_{n+p}}<\eps$, т.е. условие Коши
выполнено и для $\sum a_k$.
\end{proof}

\begin{definition}
Ряд \emph{сходится условно}, если он сходится, но не абсолютно.
\end{definition}

\section{Знакочередующиеся ряды. Признак Лейбница}

\begin{definition}
Ряд вида $\displaystyle\sum_{k=1}^\infty(-1)^{k+1}a_k=a_1-a_2+a_3-a_4+\dots$ с $a_k>0$ называют
\emph{знакочередующимся}.
\end{definition}

\begin{theorem}[Лейбниц]
Если $a_k>0$ монотонно убывают и $a_k\to0$, то знакочередующийся ряд
$\sum(-1)^{k+1}a_k$ сходится. При этом его сумма $0<S<a_1$, а остаток
$\displaystyle R_n=\sum_{k=n+1}^\infty(-1)^{k+1}a_k$ оценивается как $\abs{R_n}\le a_{n+1}$.
\end{theorem}

\begin{proof}
Чётные частичные суммы $S_{2n}=(a_1-a_2)+(a_3-a_4)+\dots+(a_{2n-1}-a_{2n})$ не убывают (каждая
скобка $\ge0$). С другой стороны, $S_{2n}=a_1-(a_2-a_3)-\dots-(a_{2n-2}-a_{2n-1})-a_{2n}<a_1$,
т.е. $\{S_{2n}\}$ ограничена сверху; по теореме Вейерштрасса $S_{2n}\to S$. Наконец
$S_{2n+1}=S_{2n}+a_{2n+1}\to S+0=S$, поэтому и $S_n\to S$. Оценка остатка следует из того, что
$R_n$ — снова лейбницевский ряд с первым членом $a_{n+1}$.
\end{proof}

\begin{example}
$\displaystyle\sum_{k=1}^\infty\frac{(-1)^{k+1}}{k}=1-\frac12+\frac13-\frac14+\dots=\ln2$ сходится
по признаку Лейбница, но \emph{условно} (ряд из модулей — гармонический, расходится).
\end{example}

\section{Перестановки членов}

\begin{theorem}[Риман]
Если ряд сходится условно, то для любого $\ell\in\R$ (в том числе $\ell=\pm\infty$) можно так
переставить его члены, что новый ряд будет сходиться к $\ell$.
\end{theorem}

\begin{theorem}[Дирихле]
Если ряд сходится абсолютно, то любая его перестановка сходится к той же сумме.
\end{theorem}

\begin{remark}
Иллюстрация теоремы Римана: переставив члены ряда
$\sum\frac{(-1)^{k+1}}{k}=S$ как
$\bigl(1-\tfrac12-\tfrac14\bigr)+\bigl(\tfrac13-\tfrac16-\tfrac18\bigr)+\dots$, получим ряд с
суммой $\tfrac12 S$ — перестановка изменила сумму, что невозможно для абсолютно сходящихся рядов.
\end{remark}
